\documentclass[conference]{IEEEtran}
\IEEEoverridecommandlockouts
% The preceding line is only needed to identify funding in the first footnote. If that is unneeded, please comment it out.
\usepackage{cite}
\usepackage{amsmath,amssymb,amsfonts}
\usepackage{algorithmic}
\usepackage{graphicx}
\usepackage{textcomp}
\usepackage{xcolor}
\usepackage{booktabs}
\usepackage{multirow}
\def\BibTeX{{\rm B\kern-.05em{\sc i\kern-.025em b}\kern-.08em
    T\kern-.1667em\lower.7ex\hbox{E}\kern-.125emX}}
\begin{document}

\title{Mini Colossal-AI: Implementing Distributed Parallelism Strategies from Scratch\\}

\author{\IEEEauthorblockN{Dinesh Chand}
\IEEEauthorblockA{\textit{MTech-AI, IISc-Bangalore} \\
dineshchand@iisc.ac.in}
\and
\IEEEauthorblockN{Goutham P}
\IEEEauthorblockA{\textit{MTech-AI, IISc-Bangalore} \\
gouthamp@iisc.ac.in}
\and
\IEEEauthorblockN{Meena Chidambaram}
\IEEEauthorblockA{\textit{MTech-AI, IISc-Bangalore} \\
meenac@iisc.ac.in}
\and
\IEEEauthorblockN{Nadhiya R S}
\IEEEauthorblockA{\textit{MTech-AI, IISc-Bangalore} \\
nadhiyar@iisc.ac.in}
}

\maketitle

\begin{abstract}
Training large language models requires distributing computation across multiple accelerators, but existing frameworks hide the complexity behind high-level APIs. In this project, we build \textit{Mini Colossal-AI}, a simplified distributed training library implemented entirely from scratch using low-level \texttt{torch.distributed} primitives. We implement four core parallelism strategies --- data parallelism with ring all-reduce, ZeRO optimizer sharding (Stage 1 and 2), 1D tensor parallelism, and pipeline parallelism with the 1F1B schedule --- and compose them into hybrid configurations through a unified plugin following Colossal-AI's 3D parallelism mesh design. We train GPT-2 on WikiText-2 across a 4-node GPU cluster and present preliminary throughput, memory, and scaling results for each strategy and their combinations.
\end{abstract}

\begin{IEEEkeywords}
Distributed Training, Data Parallelism, Pipeline Parallelism, Tensor Parallelism, ZeRO, Hybrid Parallelism
\end{IEEEkeywords}

\section{Introduction}

In Phase 1, our proposal centered on benchmarking existing PyTorch APIs (DDP, FSDP1 vs FSDP2, native pipeline) against Colossal-AI's implementations. The feedback pointed out that simply calling existing tools does not constitute sufficient project work, and that the proposal lacked a clear goal and well-defined metrics.

We have significantly revised our approach. Instead of comparing existing tools, we now \textit{build our own} distributed training library from the ground up. Every parallelism strategy is implemented using only the low-level communication primitives in \texttt{torch.distributed} --- namely \texttt{send}, \texttt{recv}, \texttt{all\_reduce}, \texttt{broadcast}, \texttt{reduce\_scatter}, and \texttt{all\_gather}. We do not call PyTorch's DDP, FSDP, or PipelineParallel APIs anywhere in our codebase.

The result is \textit{Mini Colossal-AI} (``minicolossal''), a library that replicates the core parallelism strategies offered by Colossal-AI \cite{10.1145/3605573.3605613} in a simplified, student-readable form. We train a GPT-2 model \cite{radford2019language} built from scratch on WikiText-2 \cite{merity2016pointer} across a 4-node AWS cluster, and we collect concrete metrics --- throughput, peak GPU memory, scaling efficiency, and bubble ratio --- for each strategy.

\section{Related Work}

Data Parallelism (DP) replicates the model on every GPU and synchronizes gradients after each backward pass. The main bottleneck is the all-reduce communication. Ring all-reduce \cite{patarasuk2009bandwidth} distributes the reduction across all nodes in $2(N{-}1)$ steps, achieving bandwidth-optimal communication. PyTorch's DDP \cite{li2020pytorch} further improves on this with gradient bucketing, which batches many small gradient tensors into larger buffers before triggering all-reduce, overlapping communication with the backward pass.

ZeRO \cite{rajbhandari2019zero} addresses the memory redundancy problem in DP. Standard DP stores the full optimizer states, gradients, and parameters on every GPU. ZeRO Stage 1 partitions only the optimizer states, while Stage 2 also partitions gradients, reducing per-GPU memory from $16P$ bytes to roughly $4P + 12P/N$ (Stage 1) or $4P + 4P/N$ (Stage 2), where $P$ is the parameter count and $N$ is the number of GPUs.

Tensor Parallelism (TP) splits individual layers across GPUs. Megatron-LM \cite{DBLP:journals/corr/abs-1909-08053} introduced 1D TP for transformers, where each attention head and MLP slice is computed on a separate GPU, requiring two all-reduces per transformer block. Colossal-AI extended this to 2D and 3D layouts, though these require high-bandwidth interconnects like NVLink to be practical.

Pipeline Parallelism (PP) assigns consecutive layers to different GPUs. GPipe \cite{DBLP:journals/corr/abs-1811-06965} splits mini-batches into micro-batches for concurrent execution but suffers from a pipeline flush at each step. PipeDream \cite{10.1145/3341301.3359646} introduced the 1F1B (one-forward-one-backward) schedule, which interleaves forward and backward passes to keep more stages busy simultaneously. Colossal-AI adopts the 1F1B schedule as its default pipeline strategy.

Colossal-AI \cite{10.1145/3605573.3605613} unifies these strategies through a \texttt{ProcessGroupMesh} that creates an N-dimensional grid of ranks. Each axis of the mesh corresponds to a parallelism dimension (DP, PP, TP), and the system composes them without code duplication. Our work follows this design closely.


\section{Methodology}

We implement Mini Colossal-AI as a Python library (\texttt{minicolossal/}) with a separate benchmarking suite (\texttt{benchmarks/}). All experiments run on a cluster of 4 AWS g4dn.xlarge instances, each with a single NVIDIA Tesla T4 (16\,GB VRAM), connected over a VPC private network using TCP sockets (no EFA or InfiniBand). We use PyTorch 2.4.1 with CUDA 12.1 and the NCCL backend.

\subsection{Model and Data}

We implement GPT-2 from scratch as a stack of transformer decoder blocks. For most experiments we use GPT-2 Medium (354M parameters, 24 layers, hidden size 1024, 16 attention heads). For 4-stage pipeline parallelism we use GPT-2 Large (774M parameters, 36 layers) which does not fit on a single T4 for training. The dataset is WikiText-2 (approximately 2.4M tokens), tokenized with the GPT-2 byte-pair encoding via \texttt{tiktoken}. We use a fixed sequence length of 256.

\subsection{Component 1: Data Parallelism}

Each GPU holds a full copy of the model and processes a different data shard. After the backward pass, gradients must be averaged across all GPUs before the optimizer step. We implement three gradient synchronization methods:

\begin{enumerate}
    \item \textbf{Naive all-reduce}: One \texttt{dist.all\_reduce} call per parameter tensor. Simple but slow due to many small messages.
    \item \textbf{Ring all-reduce}: We manually implement the reduce-scatter and all-gather phases using \texttt{dist.send}/\texttt{recv} in a ring topology, following the approach in \cite{patarasuk2009bandwidth}.
    \item \textbf{Bucketed all-reduce}: We batch gradient tensors into 25\,MB buckets and call \texttt{dist.all\_reduce} once per bucket. This is the same idea behind PyTorch DDP's gradient bucketing.
\end{enumerate}

In practice, the ring all-reduce with explicit \texttt{send/recv} was unstable with NCCL at 4 nodes due to point-to-point reliability issues, so our primary 4-node method uses bucketed \texttt{all\_reduce}. The ring method still works and is used for the 2-node case.

\subsection{Component 2: ZeRO Optimizer Sharding}

We implement ZeRO Stage 1 and Stage 2 as custom optimizer wrappers:

\begin{itemize}
    \item \textbf{Stage 1}: Parameters are partitioned into $N$ groups. Each GPU creates Adam optimizer states only for its $1/N$ partition. After the backward pass, gradients are all-reduced (every GPU gets the averaged gradient), each GPU runs the optimizer step on its own partition, and then updated parameters are broadcast via \texttt{all\_gather}.
    \item \textbf{Stage 2}: In addition to partitioning optimizer states, gradients are also partitioned using \texttt{reduce\_scatter}. Each GPU only receives the gradients for its own parameter partition, updates locally, then calls \texttt{all\_gather} to distribute updated parameters.
\end{itemize}

Both stages accept an optional \texttt{group} parameter so that they can operate within a sub-group of GPUs when used in hybrid parallelism configurations.

\subsection{Component 3: 1D Tensor Parallelism}

We follow the Megatron-LM \cite{DBLP:journals/corr/abs-1909-08053} approach. We implement two custom linear layers:

\begin{itemize}
    \item \textbf{ColumnParallelLinear}: Splits the weight matrix along columns. Each GPU computes a partial output; no forward communication is needed. During the backward pass, input gradients are all-reduced.
    \item \textbf{RowParallelLinear}: Splits the weight matrix along rows. Each GPU computes a partial result, then an all-reduce sums the partial outputs.
\end{itemize}

In each transformer block, the attention Q/K/V projections use column-parallel (splitting across heads), and the output projection uses row-parallel. Similarly in the MLP, the first linear is column-parallel and the second is row-parallel. This results in exactly 2 all-reduces per transformer block (1 in attention, 1 in MLP), totaling 48 all-reduces per forward pass for a 24-layer model.

We originally proposed 2D/3D tensor parallelism in Phase 1. However, these require NVLink-class bandwidth to be practical. Over our TCP network, even 1D TP shows significant communication overhead, which is itself an interesting finding. We discuss 2D/3D TP theoretically as future work.

\subsection{Component 4: Pipeline Parallelism with 1F1B}

We partition the model's transformer blocks evenly across pipeline stages (e.g., 9 blocks per stage for 36 layers on 4 GPUs). Activations are sent forward via \texttt{dist.send}/\texttt{recv}, and gradients are sent backward similarly.

We implement two pipeline schedules:
\begin{itemize}
    \item \textbf{Naive schedule}: Forward all micro-batches through stage 0, then stage 1, etc., followed by backward in reverse. Only one GPU is active at a time. Bubble ratio $\approx (S{-}1)/S$.
    \item \textbf{1F1B schedule}: After a warmup phase of $S{-}1$ forward passes, each stage alternates one forward and one backward pass in steady state, followed by a cooldown to drain remaining backwards. Bubble ratio $\approx (S{-}1)/(S{-}1+M)$, where $S$ is the number of stages and $M$ is the number of micro-batches.
\end{itemize}

With $S=4$ and $M=8$, the 1F1B schedule reduces bubble from 75\% (naive) to approximately 27\%. The 1F1B functions accept optional \texttt{prev\_rank}/\texttt{next\_rank} parameters for rank remapping in hybrid configurations.

\subsection{Hybrid Parallelism and the Unified Plugin}

Following Colossal-AI's design \cite{10.1145/3605573.3605613}, we build a unified 3D plugin rather than separate scripts for each combination. The key components are:

\begin{itemize}
    \item \textbf{ProcessGroupMesh}: An N-dimensional grid of ranks. Given a mesh shape (e.g., $\text{dp}=2, \text{pp}=2, \text{tp}=1$), it maps each global rank to coordinates and creates process sub-groups along each axis.
    \item \textbf{MiniColossalPlugin}: An orchestrator that, given \texttt{tp\_size}, \texttt{pp\_size}, and \texttt{zero\_stage}, auto-calculates $\text{dp\_size} = \text{world\_size} / (\text{tp\_size} \times \text{pp\_size})$ and composes the standalone building blocks accordingly.
\end{itemize}

The design principle is composable reuse: each standalone component exposes an optional \texttt{group} or rank-mapping parameter, so the same code runs both standalone and within hybrid configurations without duplication. Following Colossal-AI's constraint, ZeRO Stage 2 is not permitted when pipeline parallelism is active, due to the prohibitive gradient synchronization costs.

\section{Results and Discussion}

All results below are preliminary, collected on our 4-node cluster. Throughput marked with * is aggregate across all GPUs.

\subsection{Standalone Components}

Table~\ref{tab:standalone} presents the results for each component in isolation.

\begin{table}[h]
\centering
\caption{Standalone component results (4 GPUs)}
\label{tab:standalone}
\begin{tabular}{lrr}
\toprule
\textbf{Configuration} & \textbf{Throughput} & \textbf{Peak Mem} \\
\midrule
Single GPU (baseline) & 1,361 tok/s & 7.67 GB \\
DP 4-way (bucketed) & 1,322 tok/s* & 7.67 GB \\
ZeRO Stage 1 (4-way) & 926 tok/s* & 9.47 GB \\
ZeRO Stage 2 (4-way) & 873 tok/s* & 6.96 GB \\
TP 4-way & 915 tok/s & 3.39 GB \\
PP 1F1B (4 stages) & 3,040 tok/s & 2.77 GB \\
\bottomrule
\end{tabular}
\end{table}

\textbf{Data Parallelism} achieves a scaling efficiency of 24.3\% across 4 nodes. This is low, and the main reason is that our cluster uses TCP sockets over a VPC network rather than high-bandwidth interconnects like InfiniBand or EFA. Every gradient synchronization step involves sending the full model gradients ($\sim$1.3\,GB) over the network, and with only $\sim$5\,Gbps of network bandwidth, this dominates the step time. In a real data-center setting with NVLink or InfiniBand, this number would be much higher.

\textbf{ZeRO Stage 1} achieves a 4$\times$ reduction in optimizer memory (0.71\,GB across 4 GPUs vs.\ 2.83\,GB standard). The peak memory actually increases slightly (9.47 vs.\ 7.67\,GB) because the all-gather of full parameters after each step creates temporary copies. ZeRO Stage 2 further reduces memory to 6.96\,GB by also partitioning the gradients via reduce-scatter.

\textbf{Tensor Parallelism} reduces peak memory by 56\% (from 7.67 to 3.39\,GB) since each GPU holds only $1/4$ of the weight matrices. However, the 48 all-reduces per forward pass create a heavy communication bottleneck, resulting in 915 tok/s vs.\ the single-GPU 1,361 tok/s. This result highlights why real systems like Megatron-LM use tensor parallelism only within a single node (where NVLink provides high bandwidth) and data/pipeline parallelism across nodes.

\textbf{Pipeline Parallelism} with 1F1B achieves the highest throughput at 3,040 tok/s. This is because communication is minimal --- only point-to-point transfers of activation tensors between adjacent stages, rather than all-to-all gradient synchronization. Memory per stage drops to 1.59--2.77\,GB (64\% reduction). The theoretical bubble ratio with $S{=}4$ stages and $M{=}8$ micro-batches is 27.3\%.

\subsection{Hybrid Configurations}

Table~\ref{tab:hybrid} shows results when combining strategies through our unified plugin.

\begin{table}[h]
\centering
\caption{Hybrid parallelism results (4 GPUs, unified plugin)}
\label{tab:hybrid}
\begin{tabular}{lrr}
\toprule
\textbf{Configuration} & \textbf{Throughput} & \textbf{Peak Mem} \\
\midrule
DP(2) $\times$ TP(2) & 2,327 tok/s* & 4.80 GB \\
DP(2) $\times$ PP(2) & 2,281 tok/s* & 4.08 GB \\
PP(2) + ZeRO-1 in DP(2) & 1,742 tok/s* & 5.74 GB \\
\bottomrule
\end{tabular}
\end{table}

The DP$\times$TP hybrid gives a good balance: tensor parallelism halves the per-GPU model memory while data parallelism provides throughput scaling. The DP$\times$PP hybrid achieves similar throughput but with even lower memory (4.08\,GB) and a reduced bubble ratio of 11.1\% (2 stages instead of 4). Adding ZeRO-1 on top of PP further reduces optimizer memory (2$\times$ saving) at the cost of lower throughput (1,742 vs.\ 2,281 tok/s) due to the extra all-gather step after each optimizer update.

\subsection{Discussion}

A few observations stand out from these preliminary results. First, the best strategy depends heavily on the bottleneck. On our bandwidth-limited cluster, pipeline parallelism dominates because it only communicates between adjacent nodes. In contrast, DP and TP, which require all-to-all communication, suffer from the slow network. On a cluster with NVLink within nodes, the balance would shift toward TP within a node and PP across nodes, which is exactly what Colossal-AI recommends.

Second, memory and throughput are often at odds. ZeRO reduces memory but adds communication overhead. TP reduces memory but adds 48 all-reduces per forward pass. PP reduces memory and improves throughput, but introduces pipeline bubbles. The hybrid configurations attempt to find good tradeoffs, and the right choice depends on the specific hardware and model.

Third, the composable design works well in practice. Because each component accepts optional \texttt{group} and rank-mapping parameters, the same functions are called in both standalone and hybrid modes. The unified plugin adds roughly 150 lines of orchestration code on top of the existing standalone components.

\section{Conclusion}

We have implemented four core distributed parallelism strategies from scratch --- data parallelism with gradient bucketing, ZeRO Stage 1 and 2, 1D tensor parallelism, and pipeline parallelism with the 1F1B schedule --- and composed them into hybrid configurations through a unified 3D plugin modeled after Colossal-AI.

For the final phase of this project, we plan to implement a simplified auto-parallelism configurator. This would be a student-friendly version of Colossal-AI's automatic parallelism planner: given a model profile (layer count, parameter sizes, activation memory) and a hardware profile (GPU memory, measured inter-node bandwidth), it would estimate the throughput and memory for each feasible parallelism configuration and select the best one automatically. This is not meant to be a production-grade optimizer, but rather a demonstration that the ``best config'' depends on the model-hardware combination, and that automated search can find configurations that outperform manually chosen ones.

We also plan to run a more thorough set of experiments --- sweeping micro-batch counts, batch sizes, and model scales --- to better characterize the tradeoff space, and to compare our hand-built components against PyTorch's native DDP as a reference upper bound.

\newpage

% ------------------------------------------------------------------------------
% Reference and Cited Works
% ------------------------------------------------------------------------------

\bibliographystyle{ieeetr}
\bibliography{references}

% ------------------------------------------------------------------------------

\end{document}
